\chapter{Meningismus (Meningeal Signs)}

\section*{Definition}
Meningismus (meningeal signs) refers to the symptom complex of neck stiffness, headache, photophobia, and phonophobia caused by inflammation or irritation of the meninges. Classic physical examination signs include nuchal rigidity, Kernig sign, and Brudzinski sign, though their sensitivity is limited.

\section*{Pathophysiology}
Inflammation of the pia mater and arachnoid mater stimulates nociceptive nerve endings (largely by the trigeminal nerve for supratentorial meninges and the upper cervical nerves for infratentorial) producing headache and neck pain. Reflex muscle spasm of the paraspinal extensor muscles causes nuchal rigidity and resistance to passive neck flexion. Stretching of inflamed nerve roots across the spinal canal during leg extension (Kernig) or neck flexion (Brudzinski) reproduces pain and reflexive hip flexion. Increased intracranial pressure from edema or impaired CSF resorption further exacerbates symptoms.

\section*{Causes}
\begin{itemize}
  \item \textbf{Infectious meningitis:} Bacterial (\textit{N. meningitidis}, \textit{S. pneumoniae}, \textit{H. influenzae}, \textit{L. monocytogenes}), viral (enterovirus, HSV, VZV, HIV), fungal (\textit{Cryptococcus}, \textit{Histoplasma}), tuberculous
  \item \textbf{Carcinomatous meningitis:} Leptomeningeal metastasis (breast, lung, melanoma, leukemia)
  \item \textbf{Aseptic meningitis:} Drug-induced (NSAIDs, IVIg, TMP-SMX), autoimmune (lupus, sarcoidosis, IgG4-related disease), vasculitis
  \item \textbf{Subarachnoid hemorrhage:} Blood in the subarachnoid space causes chemical irritation
  \item \textbf{Parameningeal infections:} Sinusitis, mastoiditis, epidural abscess
  \item \textbf{Meningeal irritation without infection:} Post-dural puncture headache, intrathecal chemotherapy
\end{itemize}

\section*{When to See a Doctor}
See your doctor if:
\begin{itemize}
  \item Headache with fever, neck stiffness, and altered mental status (classic meningitis triad)
  \item Photophobia, phonophobia, or rash (petechiae or purpura) associated with headache and fever
  \item Rapid progression over hours
  \item Neurologic deficits (cranial nerve palsy, focal weakness, seizures) with meningeal signs
\end{itemize}

\section*{Related Symptoms}
\begin{itemize}
  \item Neck stiffness and pain with active or passive flexion
  \item Severe, diffuse headache
  \item Photophobia and phonophobia
  \item Kernig sign: resistance to knee extension with hip flexed to 90\textdegree
  \item Brudzinski sign: reflexive hip/knee flexion with passive neck flexion
  \item Tripod sign: patient leans back with hands behind to relieve meningeal stretch when sitting
  \item Altered mental status, fever, nausea/vomiting
\end{itemize}
\textbf{Important:} Classic meningeal signs (Kernig, Brudzinski) have poor sensitivity; their absence does NOT rule out bacterial meningitis -- lumbar puncture is the definitive test.

\section*{Self-Care / Home Management}
\begin{itemize}
  \item Meningismus is a medical emergency; there is no safe home management
  \item Patients with suspected meningitis should go immediately to the nearest emergency department
  \item No oral intake prior to LP if possible to avoid aspiration risk
  \item Analgesics should not delay definitive evaluation
\end{itemize}

\section*{How Your Doctor Will Evaluate You}
\begin{itemize}
  \item Lumbar puncture with CSF analysis: opening pressure, cell count with differential, glucose, protein, Gram stain, culture, viral PCR, cryptococcal antigen, cytology
  \item Blood cultures (2 sets) before antibiotics if possible, but do not delay antibiotics if LP is delayed
  \item Serum procalcitonin, CRP, CBC with differential
  \item CT brain before LP if concern for mass lesion, increased ICP, or focal neurologic deficit (to rule out herniation risk)
  \item MRI brain with contrast for meningeal enhancement in suspected carcinomatous or granulomatous meningitis
\end{itemize}

\section*{Treatment Options}
\begin{itemize}
  \item Empiric broad-spectrum antibiotics and IV dexamethasone immediately after blood cultures if bacterial meningitis is suspected (ceftriaxone + vancomycin + ampicillin)
  \item Acyclovir for suspected HSV encephalitis
  \item Antifungal therapy (amphotericin B + flucytosine) for cryptococcal meningitis
  \item Antitubercular therapy for tuberculous meningitis
  \item Supportive care: IV fluids, antipyretics, seizure prophylaxis
  \item Corticosteroids for inflammatory or autoimmune meningitis
  \item Intrathecal chemotherapy for carcinomatous meningitis
  \item Source control for parameningeal infections (sinus drainage, abscess drainage)
\end{itemize}

\clearpage
